\chapter{Orchiopexy}

\section*{Introduction \& Clinical Indications}
Orchiopexy is a surgical intervention aimed at correcting cryptorchidism, a condition where one or both testicles fail to descend into the scrotum. This procedure involves the mobilization of the undescended testicle from its abnormal location, which may be situated in the inguinal canal, abdomen, or other ectopic sites, and its repositioning into the scrotum. The surgical technique culminates in securing the testicle within a dartos pouch, preventing its re-ascent. Orchiopexy is significant not only for its cosmetic benefits but also for its critical role in preserving testicular function, facilitating cancer surveillance, and reducing the risks of complications such as testicular torsion and trauma associated with undescended testes.

\begin{itemize}
  \item Orchidopexy
  \item Testicular Fixation
  \item Undescended Testicle Repair
  \item Cryptorchidism Repair
  \item Scrotal Orchiopexy
  \item Inguinal Orchiopexy
  \item Laparoscopic Orchiopexy
\end{itemize}

The primary surgical principle of orchiopexy is to achieve a tension-free placement of the testicle into its anatomically correct position within the scrotum. Testicular descent is a complex process that typically occurs by birth or within the first six months of life, influenced by hormonal, mechanical, and genetic factors. Failure of this descent, termed cryptorchidism, results in the testicle remaining in a warmer environment, which is detrimental to germ cell development and increases the risk of malignancy.

The rationale for orchiopexy encompasses several critical objectives. First, relocating the testicle to the cooler scrotal environment enhances the potential for normal germ cell development and future fertility, particularly in unilateral cases. The optimal temperature for spermatogenesis is approximately 2-3 degrees Celsius below core body temperature. Second, a scrotal testicle is easily palpable, facilitating self-examination and clinical surveillance for testicular cancer, a risk elevated in undescended testes even post-orchiopexy. Third, securing the testicle within the scrotum significantly reduces the risk of testicular torsion and protects the testicle from external trauma. Technical goals include achieving adequate spermatic cord length without tension, preserving the delicate blood supply, and ensuring stable scrotal fixation to prevent re-ascent.

The surgical correction of undescended testicles has a rich history dating back to the 19th century, following early anatomical observations by pioneers such as John Hunter. The first documented attempts at orchiopexy occurred in the mid-1800s, but these early procedures often faced challenges due to limited understanding of testicular vascularity and surgical techniques. In 1877, Annandale performed one of the first successful orchiopexies, marking a significant advancement in the field.

The early 20th century saw further evolution of the technique. In 1899, Arthur Dean Bevan introduced the inguinal approach, emphasizing meticulous dissection of the spermatic cord to achieve sufficient length. This technique laid the groundwork for modern orchiopexy. In 1909, Torek described a method of temporarily suturing the testicle to the thigh, known as orchidopexy with external fixation, to maintain its scrotal position. This was later refined by Ombr\'edanne in the 1910s, who popularized the creation of a dartos pouch for internal fixation, a technique still widely utilized today. 

The understanding of optimal surgical timing has also evolved, shifting from later childhood to early infancy (typically 6-12 months of age) to maximize fertility potential and reduce long-term risks. The advent of laparoscopy in the late 20th century revolutionized the management of non-palpable testes, allowing for diagnostic localization and definitive repair in a minimally invasive manner.

Orchiopexy is indicated for various conditions primarily related to testicular maldescent or instability. The most common indications include:

\begin{itemize}
  \item To correct congenital cryptorchidism, where one or both testicles have not spontaneously descended into the scrotum by 6-12 months of age, following a period of observation.
  \item To address an ectopic testicle, which has descended outside the normal path, often found in the superficial inguinal pouch, perineum, or femoral region.
  \item To prevent testicular torsion, particularly in cases of previously diagnosed cryptorchidism or in patients with a "bell-clapper" deformity predisposing to torsion.
  \item To treat recurrent testicular torsion, where previous episodes indicate an anatomical predisposition requiring surgical fixation.
  \item To manage testicular re-ascent, where a previously descended testicle retracts permanently out of the scrotum, often due to a short spermatic cord or hyperactive cremasteric reflex.
  \item To facilitate easier clinical examination and self-examination for early detection of testicular malignancy, as undescended testes have a higher lifetime risk of cancer.
  \item To improve the potential for normal spermatogenesis and fertility by placing the testicle in a cooler scrotal environment.
  \item To alleviate psychological distress and cosmetic concerns associated with an empty scrotum, impacting body image and self-esteem in children and adolescents.
\end{itemize}

Orchiopexy is primarily considered a first-line surgical treatment for cryptorchidism. For infants diagnosed with an undescended testicle that has not spontaneously descended by 6-12 months of age, surgical correction is the recommended intervention. This timing is based on evidence suggesting that earlier intervention optimizes the chances of preserving testicular function and reducing long-term risks associated with maldescent.

In cases of acute testicular torsion, orchiopexy is performed as an emergency procedure, often concurrently with detorsion, to salvage the testicle and prevent recurrence. Here, it serves as a primary, urgent intervention. It can also serve as a secondary or prophylactic procedure in specific scenarios, such as managing testicular re-ascent after initial descent or as a prophylactic measure following a non-surgical detorsion of a torsed testicle where the underlying anatomical predisposition needs correction. However, its most common and critical role remains as the initial definitive treatment for cryptorchidism, aiming to prevent future complications and optimize testicular health.

\section*{Relevant Surgical Anatomy}
Understanding the surgical anatomy is crucial for the successful execution of orchiopexy. The testicles typically reside within the scrotum, which provides a cooler environment essential for spermatogenesis. Each testicle receives its blood supply from the testicular artery, a branch of the abdominal aorta, and drains via the pampiniform plexus, which forms the testicular vein. The right testicular vein drains into the inferior vena cava, while the left drains into the left renal vein. The vas deferens, responsible for transporting sperm, accompanies these vessels within the spermatic cord, along with lymphatic vessels and nerves, notably branches of the genitofemoral nerve.

The spermatic cord traverses the inguinal canal, which is an oblique passage through the anterior abdominal wall. Key anatomical landmarks include the deep inguinal ring, an opening in the transversalis fascia, and the superficial inguinal ring, an opening in the external oblique aponeurosis. The gubernaculum, a mesenchymal structure, plays a pivotal role in guiding testicular descent. In cases of cryptorchidism, the testicle may be arrested at various points along this descent path or may be located in ectopic positions such as the superficial inguinal pouch, femoral triangle, or perineum. The scrotum itself is a fibromuscular sac composed of several layers, including the dartos muscle, which is essential for creating the scrotal pouch for fixation. 

Key anatomical features include:

\begin{itemize}
  \item Testicular artery and vein
  \item Vas deferens
  \item Spermatic cord
  \item Gubernaculum
  \item Inguinal canal
  \item Dartos muscle
\end{itemize}

Meticulous dissection is required to preserve the delicate blood supply to the testicle and the vas deferens, which may be shortened or abnormally positioned in cryptorchidism.

\section*{Preoperative Workup \& Preparation}
Preoperative preparation for orchiopexy is generally straightforward, focusing on ensuring patient safety and optimizing surgical conditions.

\begin{itemize}
  \item \textbf{Medical Evaluation:} A thorough history and physical examination are performed to assess the child's overall health and confirm the diagnosis of cryptorchidism.
  
  \item \textbf{Laboratory Tests:} Routine blood tests, such as a complete blood count (CBC) and basic metabolic panel, are often performed, particularly for younger children or those with underlying conditions.
  
  \item \textbf{Imaging Studies:} For non-palpable testes, an ultrasound may be performed to locate the testicle. If it remains non-palpable, diagnostic laparoscopy is often the most accurate method to locate an intra-abdominal testicle.
  
  \item \textbf{Anesthesia Consultation:} An anesthesiologist evaluates the child to determine the safest anesthesia plan.
  
  \item \textbf{NPO Status:} Patients are instructed to be NPO for a specified period before surgery to prevent aspiration during anesthesia.
  
  \item \textbf{Medication Adjustments:} Parents are advised on which regular medications should be continued or held before surgery.
  
  \item \textbf{Informed Consent:} A detailed discussion with parents regarding the procedure, its benefits, risks, and expected postoperative course is conducted, followed by obtaining written informed consent.
  
  \item \textbf{Prophylactic Antibiotics:} A single dose of prophylactic antibiotics may be administered intravenously prior to incision, especially in cases involving higher risk of infection.
\end{itemize}

\textbf{Absolute Contraindications:}

\begin{itemize}
  \item Severe, uncontrolled systemic illness that poses an unacceptable anesthetic or surgical risk.
  \item Active infection in the inguinal or scrotal region, including cellulitis or abscess.
  \item Uncorrected coagulopathy or severe bleeding disorder that cannot be safely managed preoperatively.
\end{itemize}

\textbf{Relative Contraindications \& Precautions:}

\begin{itemize}
  \item \textbf{Prematurity/Low Birth Weight:} Surgery is often deferred in very premature infants until they are more stable.
  \item \textbf{Significant Cardiopulmonary Disease:} Requires thorough preoperative optimization and careful anesthetic management.
  \item \textbf{Testicular Agenesis:} If investigations confirm the absence of a testicle, further surgical exploration may be unnecessary.
  \item \textbf{Precautions:}
    \begin{itemize}
      \item \textbf{Vas Deferens and Testicular Vessel Preservation:} Gentle dissection is essential to avoid injury to these structures.
      \item \textbf{Tension-Free Placement:} Ensuring the testicle is placed without tension is critical to prevent ischemia.
      \item \textbf{Hemostasis:} Careful attention to hemostasis minimizes the risk of hematoma formation.
      \item \textbf{Associated Anomalies:} Awareness of potential associated anomalies is crucial for successful outcomes.
      \item \textbf{Previous Inguinal Surgery:} Prior surgery can complicate dissection and increase the risk of injury.
    \end{itemize}
\end{itemize}

\section*{Surgical Technique \& Procedure}
Orchiopexy is typically performed under general anesthesia. The patient is positioned supine, with legs slightly abducted, and the surgical area, including the groin and scrotum, is prepped with an antiseptic solution and draped in a sterile manner.

The procedure generally involves the following key steps:

\begin{itemize}
  \item \textbf{Anesthesia and Positioning:} General anesthesia is induced. The patient is placed in a supine position, and the groin and scrotum are prepped with an antiseptic solution.
  
  \item \textbf{Incision:} For palpable undescended testes, a small transverse incision (Lanz incision) is made in the inguinal crease. For non-palpable testes, diagnostic laparoscopy may precede or be combined with the open approach, often requiring a small umbilical incision for the camera port.
  
  \item \textbf{Inguinal Canal Exploration and Testicular Mobilization:} The external oblique aponeurosis is incised, and the inguinal canal is opened. The spermatic cord is identified, and the testicle is carefully dissected free from surrounding tissues. The spermatic cord is meticulously freed to achieve maximum length without tension. A patent processus vaginalis, if present, is ligated and divided.
  
  \item \textbf{Scrotal Pouch Creation:} A small incision is made in the dependent part of the scrotum, and a dartos pouch is created by dissecting a pocket between the dartos muscle and scrotal skin.
  
  \item \textbf{Testicular Transposition and Fixation:} A tunnel is created from the inguinal incision to the scrotal pouch. The mobilized testicle is guided through this tunnel into the scrotal pouch and secured within the dartos pouch using absorbable sutures to prevent re-ascent.
  
  \item \textbf{Closure:} The inguinal incision is closed in layers, and the scrotal incision is also closed in layers. No drains are typically required for routine orchiopexy.
\end{itemize}

For intra-abdominal or very high inguinal non-palpable testes, a laparoscopic approach may be employed. This can be a single-stage procedure if sufficient cord length is achieved after extensive mobilization or a two-stage Fowler-Stephens orchiopexy for very high testes requiring division of the main testicular artery.

\section*{Complications \& Risks}
While orchiopexy is generally safe, it carries potential risks categorized into general surgical complications and procedure-specific risks.

\textbf{General Surgical Risks:}

\begin{itemize}
  \item \textbf{Anesthesia Complications:} Rare adverse reactions to medications or respiratory issues.
  \item \textbf{Bleeding:} Hematoma formation can occur but is uncommon.
  \item \textbf{Infection:} Wound infection may occur, requiring antibiotics or drainage.
  \item \textbf{Pain:} Postoperative pain is expected but usually well-controlled.
\end{itemize}

\textbf{Procedure-Specific Risks:}

\begin{itemize}
  \item \textbf{Testicular Atrophy or Ischemia:} Compromise of blood supply can lead to testicular loss.
  \item \textbf{Vas Deferens Injury:} Damage can impair fertility, especially if bilateral.
  \item \textbf{Epididymal Injury:} Damage can impact fertility.
  \item \textbf{Testicular Re-ascent:} The testicle may pull back out of the scrotum over time.
  \item \textbf{Hydrocele Formation:} Fluid accumulation around the testicle may require further intervention.
  \item \textbf{Wound Complications:} Hypertrophic scarring or wound dehiscence may occur.
  \item \textbf{Nerve Injury:} Damage to sensory nerves can lead to numbness or chronic pain.
  \item \textbf{Failure to Locate Testicle:} In non-palpable cases, the testicle may not be found.
  \item \textbf{Persistent Malignancy Risk:} Increased lifetime risk of testicular cancer remains.
\end{itemize}

\section*{Postoperative Care \& Recovery}
Following orchiopexy, the child is transferred to the Post-Anesthesia Care Unit (PACU) for monitoring as they recover from anesthesia. Vital signs, pain levels, and the surgical site are continuously assessed. Most children are discharged home on the same day, as orchiopexy is commonly performed as an outpatient procedure.

During the first 24-48 hours, parents are advised to monitor for signs of complications such as excessive swelling, redness, fever, or drainage from the incision sites. Mild pain, bruising, and swelling in the groin and scrotum are normal and typically resolve within 1-2 weeks. Activity restrictions are usually in place for the first 1-2 weeks, including avoiding strenuous activities and rough play. A scrotal support or snug underwear may be recommended for comfort. 

Long-term, the goal is to ensure the testicle remains in its scrotal position and grows appropriately, with ongoing surveillance for any potential complications. Full return to normal activities, including sports, is typically permitted after 2-4 weeks, depending on the surgeon's assessment.

During orchiopexy, the surgeon documents several key findings, including the location of the undescended testicle, its size and appearance, and the integrity of the spermatic cord and vas deferens. The mobility of the testicle and the tension on the cord after mobilization are critical observations.

If a biopsy is performed, the pathology report will confirm the presence of testicular tissue and assess germ cell density. The primary pathology in cryptorchidism is the malposition itself, leading to histological changes over time, including reduced germ cell count and potential Leydig cell hyperplasia.

A normal postoperative course following orchiopexy is characterized by a swift recovery and stable positioning of the testicle within the scrotum. Patients typically experience mild to moderate pain, effectively managed with over-the-counter pain relievers. Swelling and bruising around the incision sites and in the scrotum are common and gradually subside over 1-2 weeks. 

The incision sites should heal cleanly, with absorbable sutures dissolving within 2-3 weeks. The testicle should remain palpable and securely fixed within the scrotal pouch. Children are usually able to resume light activities within a few days and return to school within a week. Full physical activity is typically permitted after 2-4 weeks, once the surgical sites are well-healed.

Postoperative care for orchiopexy focuses on wound management, pain control, and monitoring for complications.

\begin{itemize}
  \item \textbf{Wound Care:} Keep incision sites clean and dry. Sponge baths are recommended initially.
  
  \item \textbf{Pain Management:} Administer prescribed or over-the-counter pain medications as needed.
  
  \item \textbf{Activity Restrictions:} Avoid strenuous activities and contact sports for 2-4 weeks.
  
  \item \textbf{Initial Postoperative Visit:} Typically scheduled 1-2 weeks after surgery to check wound healing and assess the testicle's position.
  
  \item \textbf{Annual Physical Examinations:} Regular check-ups are recommended to monitor testicular growth and position.
  
  \item \textbf{Testicular Self-Examination Education:} Education on self-examination is vital as the child approaches adolescence.
  
  \item \textbf{Fertility Assessment:} Consider fertility assessment in adulthood for individuals with a history of bilateral cryptorchidism.
\end{itemize}

Parents and patients should be educated on warning signs such as sudden scrotal pain, increasing swelling, redness, fever, or the testicle disappearing from the scrotum, which warrant immediate medical attention.

\section*{Outcomes \& Prognosis}
Cryptorchidism is one of the most common congenital anomalies in male infants, affecting approximately 3-5\% of full-term male neonates, with a higher incidence of up to 30\% in premature infants. The condition often resolves spontaneously within the first 6 months of life, with prevalence dropping to about 1\% by one year of age. Unilateral cryptorchidism is more common than bilateral, with the right testicle affected slightly more often than the left. Approximately 80\% of undescended testes are palpable in the inguinal canal, while 20\% are non-palpable. 

Orchiopexy is one of the most frequently performed pediatric surgical procedures worldwide. The mortality rate associated with orchiopexy is exceedingly low, primarily related to general anesthesia risks, estimated to be less than 0.01\%. Morbidity rates are also low, with testicular atrophy occurring in a small percentage of cases and re-ascent being similarly uncommon.

The outcomes of orchiopexy are generally excellent, with high success rates in achieving and maintaining scrotal testicular position. Over 90-95\% of orchiopexies successfully place the testicle in the scrotum, where it remains long-term. The prognosis for fertility is significantly improved, especially in unilateral cases, where fertility rates often approach those of the general population. 

For bilateral cryptorchidism, fertility may still be reduced compared to controls, but it is substantially better than if the condition remains untreated. Early intervention is a key prognostic factor for optimizing fertility potential. Orchiopexy effectively reduces the risk of testicular torsion and protects the testicle from trauma. While the procedure facilitates early detection through self-examination, it does not completely eliminate the slightly increased lifetime risk of testicular malignancy associated with cryptorchidism. Prognostic factors influencing outcomes include the initial position of the testicle, the age at which surgery is performed, and the presence of associated anomalies.

\section*{References}
\begin{enumerate}
  \item MedlinePlus. Orchiopexy. U.S. National Library of Medicine; 2024. Available from: \url{https://medlineplus.gov/ency/article/002964.htm}
  
  \item Wein AJ, Kavoussi LR, Partin AW, Peters CA, editors. Campbell-Walsh-Wein Urology. 12th ed. Philadelphia, PA: Elsevier; 2021.
  
  \item American Urological Association. AUA Guideline: Cryptorchidism. 2014. Available from: \url{https://www.auanet.org/guidelines-and-quality/guidelines/cryptorchidism-guideline}
  
  \item Holcomb GW, Murphy JP, St. Peter SD, editors. Holcomb and Ashcraft's Pediatric Surgery. 7th ed. Philadelphia, PA: Elsevier; 2020.
\end{enumerate}

\clearpage
